\section{学习会收敛到均衡吗}
\label{sec:17-Constrained-Guarantees/06-Games-and-Equilibria/04}

一个 GAN 训练了三天。判别器的准确率在 \(0.5\) 到 \(0.8\) 之间来回摆，生成样本的质量时好时坏，没有任何一个指标单调地朝某个方向走。把学习率调小一半，摆动周期变长，摆幅没变。把生成器权重做指数滑动平均之后取样，图像质量明显好过任何一个单独的检查点。

这三条现象——不收敛的轨迹、不随步长消失的摆动、以及平均之后反而更好——不是调参没调好。它们是前面两个单元的结论合在一起之后必然出现的样子。

上一个单元造了一个对世界零假设的承诺：遗憾增长得比 \(T\) 慢。这一个单元造了一个静态的解概念：谁也无法单方面改进。这一节把两者接上——如果博弈的每一方都在跑一个无遗憾算法，长期会走到哪里。

\subsection{零和博弈里时间平均收敛到均衡}

设两人零和博弈的收益矩阵是 \(A\)，行方要最大化、列方要最小化。第 \(t\) 轮双方分别给出混合策略 \(p_t\) 与 \(q_t\)，行方承受的``损失''是 \(-p^{\trans}Aq_t\) 那一族线性函数，列方承受的是 \(p_t^{\trans}Aq\)。两边各自跑一个无遗憾算法——比如\hyperref[sec:17-Constrained-Guarantees/04-Online-Learning-and-Regret/02]{``无遗憾是可以做到的''}里的指数权重，各自把对手当作一串不需要解释的损失函数。

记时间平均 \(\bar p=\frac1T\sum_t p_t\)，\(\bar q=\frac1T\sum_t q_t\)，以及实际发生的平均收益 \(\bar u=\frac1T\sum_t p_t^{\trans}Aq_t\)。

\begin{theorem}
若行方的遗憾不超过 \(R^1_T\)、列方的遗憾不超过 \(R^2_T\)，则 \((\bar p,\bar q)\) 是这个博弈的 \(\varepsilon\)-均衡，其中 \(\varepsilon=(R^1_T+R^2_T)/T\)；并且 \(\abs{\bar u-v}\le\varepsilon\)，\(v\) 是博弈的值。
\end{theorem}

推导几行就完。把两边的遗憾定义写出来，除以 \(T\)，用效用关于各自变量的线性性把时间平均搬进去：

\[
\begin{aligned}
\max_{p}\;p^{\trans}A\bar q&\;\le\;\bar u+\frac{R^1_T}{T},\\
\min_{q}\;\bar p^{\trans}Aq&\;\ge\;\bar u-\frac{R^2_T}{T}.
\end{aligned}
\]

第一行说的是：行方对着 \(\bar q\) 打出的最好成绩，不会比它实际拿到的平均收益高出多少，否则它的遗憾就超标了。第二行对称。

再用\hyperref[sec:17-Constrained-Guarantees/06-Games-and-Equilibria/01]{``最坏情况下的最优与交换次序''}里的极小极大定理夹逼。一方面 \(v=\max_p\min_q\ge\min_q\bar p^{\trans}Aq\ge\bar u-R^2_T/T\)；另一方面 \(v=\min_q\max_p\le\max_p p^{\trans}A\bar q\le\bar u+R^1_T/T\)。于是 \(\bar u\) 被夹在 \(v\) 的两侧 \(\varepsilon\) 之内，代回上面两式即得 \((\bar p,\bar q)\) 的近似均衡性质。

代入 Hedge 的界 \(R_T=O(\sqrt{T\ln N})\)，得 \(\varepsilon=O(\sqrt{\ln N/T})\)，随 \(T\) 归零。

这个结论的分量值得掂一下。双方都没有解方程，没有算不动点，甚至没有假设对面是理性的——每一方都只是把对手当成一串任意的损失函数，跑一个对最坏情况有保证的算法。而两个各自谨慎的算法互相对打，居然把博弈的值算了出来。

顺带得到一件事：上面的推导若不引用极小极大定理，而是直接由两式相加，得到 \(\max_p\min_q\ge\min_q\max_p-\varepsilon\)，令 \(T\to\infty\) 便有 \(\max\min\ge\min\max\)，与那条永远成立的反向不等式合起来正是极小极大定理。也就是说，无遗憾动力学给了极小极大定理一个构造性证明，全程不用 Brouwer 或 Kakutani。\hyperref[sec:17-Constrained-Guarantees/06-Games-and-Equilibria/02]{``Nash 均衡存在但可能不唯一也不好''}里说 Kakutani 只给存在性、不给找到它的办法——在零和这个特例上，办法在这里。

\subsection{收敛的是时间平均，不是当前策略}

定理里出现的是 \(\bar p\) 和 \(\bar q\)，不是 \(p_T\) 和 \(q_T\)。这个区别不是证明技术的遗留，它是实质性的。

石头剪刀布是最干净的反例。双方都跑指数权重，谁也不会稳定下来：对方出石头多了，我就往布上加权重；我布多了，对方就往剪刀上加权重。策略在单纯形里绕着均匀分布转圈。连续时间下有一个确切的结论——乘性权重动力学保持``当前策略到均衡的 KL 散度''不变，于是轨迹是一条闭合的轨道，永远不收缩。离散时间下步长固定时更糟，轨迹是向外的螺旋，逐渐贴到单纯形的边界上。这两条都是定理，不是模拟观察。

而时间平均照样收敛，因为绕圈的轨迹平均下来正好落在圆心。

\begin{center}
\begin{tikzpicture}[x=1cm,y=1cm,font=\small,>=stealth]
  \draw[->,black!60] (0,0) -- (4.8,0) node[right,font=\footnotesize] {\(p\)};
  \draw[->,black!60] (0,0) -- (0,4.8) node[above,font=\footnotesize] {\(q\)};
  \begin{scope}[shift={(2.2,2.2)}]
    \draw[blue!65!black,thick,samples=240,domain=0:18.9,smooth]
      plot ({(0.55+0.045*\x)*cos(\x r)},{(0.55+0.045*\x)*sin(\x r)});
    \draw[orange!75!black,thick,samples=200,domain=1.2:18.9,smooth]
      plot ({(1.6/\x)*cos(\x r)},{(1.6/\x)*sin(\x r)});
    \fill[black] (0,0) circle (2.2pt);
  \end{scope}
  \draw[blue!65!black,thick] (5.0,3.55) -- (5.45,3.55);
  \node[anchor=west,align=left,font=\footnotesize] at (5.52,3.55)
    {当前策略：绕着均衡转，\\ 半径不收缩甚至变大};
  \draw[orange!75!black,thick] (5.0,2.15) -- (5.45,2.15);
  \node[anchor=west,align=left,font=\footnotesize] at (5.52,2.15)
    {时间平均：收敛到均衡};
  \fill[black] (5.22,1.35) circle (2.2pt);
  \node[anchor=west,align=left,font=\footnotesize] at (5.52,1.35)
    {黑点：博弈的均衡};
\end{tikzpicture}

\emph{两条曲线来自同一次运行。蓝色是每一轮的联合策略，橙色是它到该轮为止的平均。定理管的是橙色那条，日志里看到的是蓝色那条。}
\end{center}

图里蓝色那条永远到不了中心点，橙色那条却稳稳地走进去。训练日志、监控面板、单个检查点看到的都是蓝色那条，所以``没有收敛''这个观感是对的——它只是没有在定理承诺的那个量上做出判断。

一般和博弈里还要退一步。无遗憾动力学的经验联合分布收敛到的是\emph{粗相关均衡}，而不是 Nash。粗相关均衡是联合动作上的一个分布，它要求的是：在不知道这一轮抽到什么的前提下，没有人愿意事先承诺换一个固定动作。这比 Nash 弱两层——Nash 要求分布是各方边缘分布的乘积（互相独立），粗相关均衡不要求；Nash 允许偏离者在看到自己被推荐的动作之后再决定，粗相关均衡不允许。介于两者之间的是相关均衡，把``遗憾''换成更强的交换遗憾就能收敛到它。

这个层级不是理论上的挑剔。想收敛到 Nash 而不只是粗相关均衡，必须付出更多，而上一节提到的 PPAD-完全性说明这份代价不太可能是多项式的：若存在一套高效的、通用的动力学收敛到 Nash，一整类被认为难解的问题就跟着变简单了（\hyperref[sec:09-Convex-Optimization/09-Complexity-and-Approximation/01]{``P、NP 与 NP-hard：给问题按计算难度分类''}）。所以``学习会不会收敛到 Nash''这个问题，在一般和博弈里的答案大概率是否定的，而且理由是计算性的而不是动力学性的。

零和之所以是例外，是因为极小极大定理把它压成了一个凸问题：那里的粗相关均衡的边缘分布恰好就是均衡策略，弱概念自动升级成了强概念。丢掉零和，这条升级通道就断了。

\subsection{GAN 为什么摆动}

回到开头那三天的训练。GAN 的目标是一个极小极大问题：

\[
\min_G\max_D V(G,D),
\]

生成器要最小化，判别器要最大化。\hyperref[sec:14-Deep-Learning/07-Generative-Models-Unified/03]{``GAN 把散度的估计交给判别器''}讲的是这个 \(\max_D\) 在做什么——它在估计一个散度，估得越准，外层的 \(\min_G\) 才越接近\hyperref[sec:14-Deep-Learning/07-Generative-Models-Unified/01]{``四条路线在最小化哪个散度''}里那条散度极小化的路线。

而训练是怎么跑的？交替若干步梯度上升与梯度下降。这正是两个玩家各自在做近似的最优反应，与本节前半段的设定同构。于是上面每一条结论都直接落地：

轨迹不收敛是预期之内的。即使把问题理想化成完全的凸凹情形——比如 \(\min_x\max_y xy\) 这个最简单的双线性例子——固定步长的同步梯度上升下降也会向外发散，因为更新矩阵的特征值在单位圆外。步长调小只是让发散变慢，不改变方向。这是定理，不是经验。判别器准确率的摆动与它是同一件事。

平均之后更好也是预期之内的。定理保证的是时间平均，而生成器权重的指数滑动平均正是在做这件事的近似。这里必须说清它\emph{只是}近似：定理里的平均是策略（分布）的平均，而权重平均是参数的平均，两者只在参数到策略的映射是线性时才一致，神经网络远不满足这一条。所以``EMA 的样本更好''是一条被反复观察到的经验规律，理论提供的是它为什么值得一试，不是它一定管用的保证。

其余的稳定化手段可以按同样的标准分档。乐观梯度、外推梯度这一类在凸凹情形下确实有最后一轮收敛的证明，把它们搬到非凸非凹的网络上就只剩启发意义。谱归一化、梯度惩罚、双时间尺度更新则从一开始就是经验规律：它们试图让判别器更接近``那一侧真正的最优反应''，从而让上面那套分析的前提更接近成立，但没有定理保证在深度网络上做到了这一点。分不清这两档，就会把一次调参失败误解成理论错了。

自对弈落在同一个框架里。两人零和的完美信息博弈中，自对弈加上无遗憾更新有理论支撑；扑克类不完全信息博弈里的反事实遗憾极小化更是逐字的无遗憾算法，它收敛到 Nash 的理由就是本节开头那个定理。但一旦训练目标不是严格零和——多方博弈、带辅助奖励的智能体、人机协作场景——保证就退回到粗相关均衡那一档。\hyperref[sec:16-Control-and-Reinforcement-Learning/02-Reinforcement-Learning/01]{``MDP 把决策问题写成随机过程''}里那个单智能体的框架在这里也不再够用：环境里坐着另一个在学习的人，转移概率随对方的策略而变，MDP 的平稳性前提不成立。

\subsection{这一部分收在哪里}

把本单元的四节连起来看，得到的是一串逐级退让的结论。零和、混合策略、有限选项之下，博弈有唯一的值，两个操作可以交换，无遗憾学习能把它算出来。丢掉零和，值这个概念消失，均衡还在但可能不唯一，选哪一个需要数学之外的假设。均衡即使唯一，也可能对所有人都更差，而``差多少''有上界但上界依赖延迟函数的形状。想改善，能动的是规则而不是参与者，可改规则本身撞上另一组不可能定理。最后，即使均衡存在且我们想要它，学习动力学收敛的也只是时间平均，一般和情形下连收敛目标都降了一级。

每一步退让都有确切的理由，而每个理由都是同一种形状：想同时要几件好事，某条数学结论说不行。这条主线贯穿整个部分——想要同时保证隐私与精度，想要同时满足几条公平性定义，想要同时保证鲁棒与准确，想要同时让均衡稳定且高效。这些结论没有一条是暂时的技术限制，它们都是可以证明的。

面对这类结论，能做的事情是明确的：把``全都要''换成``要哪一个、放弃哪一个、放弃到什么程度''，然后把这个选择写下来，让它接受审视。这一部分给出的每一条界，都是为了让这句话能被写得足够具体。
